\scanpage{93}
\mspage{86}
% archive sheet number: 92
\[
X_{i_J}=\varinjlim_{i\in J}X_i.
\]
De tels syst. inductifs seront dits ``\uline{parfaitement adaptés}'' (à la filtration cardinale donnée, pour le cardinal $\pi$ donné --- $\pi_0$ n'y figure plus, si ce n'est comme début de la filtr. cardinale en \unsure{place|plus}.).

Comme pratiquement, on aura à considérer des $E$ stables par \inserted{petites} $\varinjlim$ filtrantes, et des filtrations cardinales stables \deleted{\unsure{par}} \unsure{et coniques|et canoniques} (renforcement de la condition b) des filtr. cardinales), c'est le cas, \deleted{\unsure{les}} et la notion des représentations \unsure{parfaitement|bien} adaptées de $X$ comme $\varinjlim_I\unsure{X_i|\alpha_i}$ ($X_i\in\mathrm{Filt}^{\pi}$), que je mettrai en avant. Le cas ``général'', où on suppose simplement $E$ satisfait $L_{\pi_0}$, est plus alambiqué et se traite moins joliment.
% slanted
\margin{On parlera de filtration cardinale \uline{stricte}, si les \uline{$\mathrm{Filt}^{\pi}$} sont stables par $\varinjlim$ filtrantes de cardinal $\leq\pi$. (et $E$ par les limites \unsure{filtrantes} petites)}

\uline{Définition 2} Soit $E$ muni d'une filtration cardinale stricte $(\underline{\mathrm{Filt}^{\pi}})_{\pi\geq\pi_0}$. \deleted{Une} \inserted{Soit $M$ une} partie \deleted{\unsure{$M$}} de $\operatorname{Ob}E$, $\underline M$ la sous-cat. pleine de $E$ qu'elle définit, \inserted{et soit $\pi$ un cardinal \deleted{\unsure{$\pi$}} $\geq\pi_0$.} On dit que \unsure{$M$} est \uline{$\pi$-accessible}, ou que $\underline M$ est $\pi$-accessible,
